% ============================================================
%  以異質圖神經網路預測短影音平台表現
%  系統設計、模型建構與實驗分析
%  概念驗證技術報告（第二版）
%
%  編譯：xelatex pipeline_report.tex（需跑兩次更新目錄）
%  Ubuntu 安裝：sudo apt-get install texlive-full texlive-xetex fonts-noto-cjk
% ============================================================
\documentclass[12pt, a4paper]{article}
\usepackage[UTF8, scheme=plain]{ctex}
\usepackage[top=2.5cm, bottom=2.5cm, left=3cm, right=3cm]{geometry}
\usepackage{setspace}\onehalfspacing
\usepackage{amsmath, amssymb}
\usepackage{bm}
\usepackage{tikz}
\usetikzlibrary{arrows.meta, positioning, shapes.geometric, shapes.arrows,
                fit, backgrounds, calc, decorations.pathreplacing, matrix}
\usepackage{booktabs, multirow, tabularx, array}
\usepackage{xcolor}
\usepackage{listings}
\usepackage[colorlinks=true, linkcolor=blue!60!black,
            citecolor=green!50!black, urlcolor=blue!60!black]{hyperref}
\usepackage{natbib}
\usepackage{float, caption, subcaption, graphicx}
\usepackage{enumitem}
\usepackage{tcolorbox}
\tcbuselibrary{skins, breakable, listings}
\usepackage{fancyhdr}
\pagestyle{fancy}\fancyhf{}
\rhead{\small 短影音平台表現預測：技術報告 v2}
\lhead{\small \leftmark}\cfoot{\thepage}

%% ── 顏色 ────────────────────────────────────────────
\definecolor{vidblue}{RGB}{41, 128, 185}
\definecolor{authgreen}{RGB}{39, 174, 96}
\definecolor{catred}{RGB}{192, 57, 43}
\definecolor{lightbg}{RGB}{248, 249, 250}
\definecolor{codebg}{RGB}{245, 245, 245}
\definecolor{hlgold}{RGB}{241, 196, 15}
\definecolor{darktext}{RGB}{44, 62, 80}
\definecolor{tipsbox}{RGB}{232, 245, 233}
\definecolor{warnbox}{RGB}{255, 243, 205}
\definecolor{infobox}{RGB}{227, 242, 253}

%% ── tcolorbox 樣式 ──────────────────────────────────
\newtcolorbox{tipbox}[1][]{
  enhanced, colback=tipsbox, colframe=authgreen!70, fonttitle=\bfseries,
  title={#1}, arc=4pt, left=4pt, right=4pt}
\newtcolorbox{warnbox2}[1][]{
  enhanced, colback=warnbox, colframe=hlgold!80, fonttitle=\bfseries,
  title={#1}, arc=4pt, left=4pt, right=4pt}
\newtcolorbox{infobox2}[1][]{
  enhanced, colback=infobox, colframe=vidblue!70, fonttitle=\bfseries,
  title={#1}, arc=4pt, left=4pt, right=4pt}
\newtcolorbox{analogybox}[1][]{
  enhanced, colback=gray!6, colframe=gray!40, fonttitle=\bfseries\itshape,
  title={#1}, arc=4pt, breakable}

%% ── 程式碼樣式 ──────────────────────────────────────
\lstset{
  backgroundcolor=\color{codebg}, basicstyle=\ttfamily\small,
  breaklines=true, frame=single, rulecolor=\color{gray!40},
  numbers=left, numberstyle=\tiny\color{gray}, xleftmargin=1.5em,
  language=bash, keywordstyle=\color{blue!70!black}\bfseries,
  commentstyle=\color{gray}\itshape}

%% ── TikZ 樣式 ───────────────────────────────────────
\tikzset{
  vidnode/.style  ={draw=vidblue,  fill=vidblue!15, rounded corners=4pt,
                    minimum width=1.8cm, minimum height=0.7cm,
                    font=\small\bfseries, text=darktext},
  authnode/.style ={draw=authgreen,fill=authgreen!15, rectangle,
                    minimum width=1.8cm, minimum height=0.7cm,
                    font=\small\bfseries, text=darktext},
  catnode/.style  ={draw=catred,   fill=catred!15,   diamond, aspect=2.2,
                    minimum width=2.4cm, minimum height=0.7cm,
                    font=\small\bfseries, text=darktext},
  layerbox/.style ={draw=gray!60,  fill=lightbg, rounded corners=6pt,
                    minimum width=7cm, minimum height=1.0cm,
                    font=\small, text=darktext, align=center},
  arrowstyle/.style={-Stealth, thick, draw=gray!70},
}

%% ── 標題 ────────────────────────────────────────────
\title{
  \vspace{-1cm}
  {\LARGE\bfseries 以異質圖神經網路預測短影音平台表現}\\[0.5em]
  {\large 系統設計、模型建構與實驗分析}\\[0.4em]
  {\normalsize\itshape 概念驗證技術報告（第二版）}
}
\author{roy891030}
\date{2026 年 5 月}

\begin{document}
\maketitle\thispagestyle{fancy}

% ────────────────────────────────────────────────────
\begin{abstract}
短影音平台每天上傳數十萬支新影片，如何在新影片上架的第一時間估計它的互動潛力，
是平台流量調度的核心問題。傳統的推薦系統需要「使用者歷史互動資料」才能運作，
對一支剛上傳、沒有任何人看過的新影片完全束手無策。

本報告設計了一套三管線研究框架，其中核心創新（Pipeline C）建立了一張以影片為
中心的關係網路（異質圖），把影片的視覺內容、作者背景、同分類影片三種信號用
數學方式傳遞融合，讓一支全新影片也能「借用鄰居的資訊」來估計自己的表現。

在 RTX~4090 GPU 上訓練結果：平均觀看時長的排名相關係數達 \textbf{0.256}，
比不使用圖結構的傳統方法高出 4 倍以上。
\end{abstract}

\tableofcontents\newpage

% ════════════════════════════════════════════════════
\section{讀前必讀：關鍵概念白話速覽}
\label{sec:concepts}

本節幫助沒有深度學習背景的讀者快速理解後文會用到的概念，
有相關基礎的讀者可以跳過。

\subsection*{什麼是「特徵向量」？}

電腦無法直接理解「影片的畫面很有趣」，它只認識數字。
\textbf{特徵向量}就是用一串數字來代表一個物件的特性。
例如把一支影片的畫面送進 ViT-B/16（一種圖像理解 AI 模型）後，
會輸出一個 768 個數字組成的清單——這 768 個數字共同描述了這支影片的視覺風格，
就像一支影片的「視覺指紋」。

\subsection*{什麼是「圖」？}

這裡的圖不是長條圖，而是數學上的\textbf{圖論（Graph Theory）}概念：
\begin{itemize}
\item \textbf{節點（Node）}：圖中的每一個「點」，代表一個實體。例如每支影片是一個節點。
\item \textbf{邊（Edge）}：連接兩個節點的「線」，代表兩者之間的關係。例如「這支影片是這個作者拍的」就是一條邊。
\end{itemize}
社群網路（Facebook 的好友關係）、交通網路（城市之道路）都是圖的例子。

\subsection*{什麼是「圖神經網路（GNN）」？}

普通的神經網路只能處理固定格式的表格資料。
GNN 則是設計來處理「有關係的資料（圖）」的神經網路。

\begin{analogybox}[生活中的類比：問問鄰居]
想像你要猜一首新歌會不會流行，但你對這首歌一無所知。
你可以問：這個歌手過去的歌都怎樣？這個曲風類型最近有沒有在流行？
和這首歌最像的其他歌有沒有爆紅？

GNN 做的事情完全一樣：讓每個節點「問問自己的鄰居」，
把鄰居的資訊收集過來，更新自己的表示向量。
對一支新影片來說，它的鄰居是作者、所屬分類、以及視覺風格最相似的其他影片。
\end{analogybox}

\subsection*{什麼是「訊息傳遞」？}

GNN 的核心運作機制叫做\textbf{訊息傳遞（Message Passing）}。
每一輪（稱為一層）的過程如下：
\begin{enumerate}
\item 每個節點把自己目前的特徵向量傳給所有相連的鄰居（「傳遞訊息」）
\item 每個節點收集來自鄰居的訊息，用數學方式聚合（通常是加總或平均）
\item 每個節點用聚合後的結果更新自己的特徵向量
\end{enumerate}
跑 2 層 GNN，每個節點就能「感知到」距離它 2 步以內的所有鄰居的資訊。

\subsection*{什麼是「冷啟動」？}

推薦系統最頭痛的問題之一。一支全新影片上傳後，沒有任何人看過它，
所以沒有任何互動資料（按讚、觀看時長等）。這種「一片空白」的狀態就叫\textbf{冷啟動（Cold Start）}。
傳統的協同過濾系統（根據「你看過的東西和其他人看過的東西」來推薦）在冷啟動時完全失效。
本報告的 Pipeline C 刻意設計成：即使沒有任何互動記錄，也能用影片的視覺內容和作者背景做預測。

\subsection*{什麼是「Spearman 相關係數」？}

\textbf{Spearman 相關係數}衡量「預測排名」和「真實排名」的一致程度：
\begin{itemize}
\item $= +1.0$：預測排名和真實排名完全一致（最好）
\item $= 0.0$：和瞎猜一樣，沒有相關
\item $= -1.0$：完全相反（預測第 1 名的，真實上排最後）
\end{itemize}
在推薦系統中，排名能力比絕對數值更重要——我們更在意「哪支影片應該排第一」，而不是「按讚率精確到幾位小數」。

% ════════════════════════════════════════════════════
\section{研究動機：為什麼要做這件事？}
\label{sec:motivation}

\subsection{短影音平台的三個核心痛點}

\begin{enumerate}[leftmargin=*, label=\textbf{痛點 \arabic*.}]

\item \textbf{新影片冷啟動無解}

  平台每天上傳幾十萬支新影片。在最初幾小時，這些影片沒有任何歷史互動資料。
  主流推薦演算法（協同過濾）完全依賴「這個使用者過去喜歡過什麼」和
  「和你口味相似的人喜歡什麼」——新影片兩者皆無，只能靠運氣分配初始流量。

\item \textbf{互動行為的極度偏斜}

  並非所有影片都能獲得均等的互動。表~\ref{tab:label_stats} 顯示，
  按讚率的中位數是 \textbf{0}——意思是超過一半的影片沒有人按讚。
  這種「贏者全拿」的分佈讓任何預測模型都面臨嚴峻挑戰。

\item \textbf{影片之間的語義關聯被忽視}

  傳統機器學習模型（如梯度提升樹）只能看每支影片自己的特徵，
  無法利用「同一個作者的其他影片也很受歡迎」或「和這支影片視覺風格最像的 10 支影片都爆紅了」
  這類\textbf{圖結構信號}。這正是 GNN 能夠補充的地方。

\end{enumerate}

\subsection{本報告的核心研究問題}

\begin{infobox2}[核心問題]
在不依賴任何使用者互動歷史的前提下，能否透過建立以影片為中心的關係圖，
讓圖神經網路的訊息傳遞機制有效地融合「作者背景」、「分類信號」和
「相似影片的歷史表現」，使冷啟動場景下的影片表現預測優於純特徵的傳統方法？
\end{infobox2}

% ════════════════════════════════════════════════════
\section{資料集與取得方式}
\label{sec:data}

\subsection{ShortVideo 資料集介紹}

本研究使用清華大學 \citet{shang2025shortvideo} 發表於 WWW 2025 的 ShortVideo 資料集，
從某真實短影音平台採集，是目前公開規模最大、行為類型最完整的短影音研究資料集之一。

\begin{table}[H]
\centering
\caption{資料集基本統計（本報告使用版本）}
\label{tab:dataset}
\begin{tabular}{lrl}
\toprule
\textbf{項目} & \textbf{數量} & \textbf{說明} \\
\midrule
曝光記錄總數   & 6,767,010 & 每列 = 一位使用者觀看一支影片一次 \\
獨立使用者數   & 10,000 & 已雜湊匿名 \\
獨立影片數     & 153,561 & 含視覺特徵的完整集合 \\
資料時間跨度   & 7 天 & 2022-09-16 至 2022-09-22 \\
\midrule
有效標籤影片   & 98,999 & 曝光 $\geq 10$ 次，標籤統計較穩定 \\
Pipeline C 工作集 & \textbf{54,088} & 有標籤 $\cap$ 有預抽特徵的影片 \\
\bottomrule
\end{tabular}
\end{table}

\begin{warnbox2}[「Tiny」的名字會誤導人]
這份資料集雖然叫「tiny version」，但它包含了完整的 7 天互動日誌（670 萬行）
和所有 153,561 支影片的預抽特徵。「Tiny」指的是原始 .mp4 影片檔只附了 10 支示範，
\textbf{互動資料和視覺特徵是完整的}，對 PoC 完全足夠。
\end{warnbox2}

\subsection{為什麼工作集只有 54,088 部？}

Pipeline C 需要兩樣東西：（1）影片的行為標籤（按讚率等），（2）影片的視覺特徵向量。
兩者來自不同的資料檔案，覆蓋的影片集合有差異：

\begin{center}
\begin{tikzpicture}
\draw[fill=vidblue!15, draw=vidblue] (0,0) ellipse (3.0cm and 1.2cm);
\draw[fill=authgreen!15, draw=authgreen] (1.5,0) ellipse (2.5cm and 1.0cm);
\node[font=\small] at (-1.5, 0) {有行為標籤};
\node[font=\small] at (-1.5,-0.4) {98,999 部};
\node[font=\small] at (3.3, 0) {有視覺特徵};
\node[font=\small] at (3.3,-0.4) {65,030 部};
\node[font=\small\bfseries] at (1.2, 0) {工作集};
\node[font=\small] at (1.2,-0.4) {54,088 部};
\end{tikzpicture}
\end{center}

另外 44,911 部有標籤但沒有特徵的影片，需要從清華主站下載 \texttt{video\_feature\_total/}
後才能納入訓練。

\subsection{行為標籤的分佈特性}

\begin{table}[H]
\centering
\caption{8 個行為指標的分佈（54,088 部影片工作集）}
\label{tab:label_stats}
\begin{tabular}{lrrrr}
\toprule
\textbf{指標} & \textbf{中位數} & \textbf{均值} & \textbf{零值比例} & \textbf{預測難度} \\
\midrule
\texttt{like\_rate}            & 0.000 & 0.026 & 84.6\% & 極難 \\
\texttt{comment\_rate}         & 0.000 & 0.004 & 96.7\% & 極難 \\
\texttt{follow\_rate}          & 0.000 & 0.004 & 96.9\% & 極難 \\
\texttt{collect\_rate}         & 0.000 & 0.006 & 95.2\% & 極難 \\
\texttt{forward\_rate}         & 0.000 & 0.002 & 98.0\% & 極難 \\
\texttt{hate\_rate}            & 0.000 & 0.001 & 99.2\% & 極難 \\
\texttt{effective\_view\_rate} & 0.833 & 0.791 &  1.8\% & 中等 \\
\texttt{mean\_watch\_time}     & 22 秒 & 33 秒 &  0.2\% & \textbf{最容易} \\
\bottomrule
\end{tabular}
\end{table}

\textbf{為什麼「rate」類指標那麼難預測？}
因為大部分影片的按讚率都是 0——你預測所有影片都是 0，準確率高達 84.6\%，
但這毫無意義，因為你完全沒有區分「爆紅的 15.4\%」和「路人的 84.6\%」的能力。
只有 \texttt{mean\_watch\_time}（觀看時長）的分佈相對正常，
且最直接受影片內容品質影響，是最值得用 GNN 預測的指標。

\subsection{資料取得完整指南（SSH 環境）}

\subsubsection*{資料來源總覽}

\begin{table}[H]
\centering
\caption{三個資料來源的功能與位置}
\begin{tabular}{lll}
\toprule
\textbf{來源} & \textbf{內容} & \textbf{位置} \\
\midrule
Dropbox A（行為日誌） & interaction.csv + 10 示範影片 & 見下方連結 \\
Dropbox B（推薦特徵） & image\_feat.npy + text\_feat.npy & 見下方連結 \\
清華主站（完整版）    & 全部影片 .npy + 文字 + 原始影片 & 帳密見下方 \\
\bottomrule
\end{tabular}
\end{table}

\subsubsection*{下載方法（Ubuntu SSH，無需 sudo）}

\begin{lstlisting}[language=bash, caption=推薦預處理特徵（Pipeline C 必要）]
cd ~/shortvideo_mmgcn_poc
# Dropbox B：image_feat.npy + text_feat.npy（約 500 MB）
wget -O /tmp/rec.zip \
  "https://www.dropbox.com/scl/fo/ha0e0wolgqgg5qskr52l1/AHZUySejwWJzfyJy8WGo2k4\
?rlkey=xwpqosx7b906yb7g8nwy53oqh&st=slry1d1l&dl=1"
unzip /tmp/rec.zip -d data_raw/video_rec_dataset/
\end{lstlisting}

\begin{lstlisting}[language=bash, caption=行為日誌 Tiny 版（Pipeline B/C 必要）]
# Dropbox A：interaction.csv + 示範影片（約 200 MB）
wget -O /tmp/tiny.zip \
  "https://www.dropbox.com/scl/fo/5z7pwp6xjkrr1926vreu6/AFYter5C6BDTOCpxkxF0k9Y\
?rlkey=p28j6u1fl1ubb7bufiq16onbl&st=7aktff85&dl=1"
unzip /tmp/tiny.zip -d data_raw/shortvideo_tiny/
\end{lstlisting}

\begin{lstlisting}[language=bash, caption=清華主站（完整版，需帳密）]
# 帳號：videodata   密碼：ShortVideo@10000
# 下載互動主檔（.csv 無連線數限制）
wget --user=videodata --password='ShortVideo@10000' -c \
     -O data_raw/shortvideo_full/interaciton_filtered.csv \
     "http://fi.ee.tsinghua.edu.cn/datasets/short-video-dataset/interaciton_filtered.csv"

# 下載 SBERT 用的標題文字（.txt 無連線數限制）
wget --user=videodata --password='ShortVideo@10000' -r -np -A "*.txt" \
     -P data_raw/shortvideo_full/title_en/ \
     "http://fi.ee.tsinghua.edu.cn/datasets/short-video-dataset/title_en/"

# 大型 .npy 檔案（限制 3 個同時連線）
wget --user=videodata --password='ShortVideo@10000' -c \
     -O data_raw/shortvideo_full/image_feat.npy \
     "http://fi.ee.tsinghua.edu.cn/datasets/short-video-dataset/image_feat.npy"
\end{lstlisting}

\begin{tipbox}[最快提升訓練效果的方式（不需要下載）]
\texttt{interaction.csv} 裡已有中文 \texttt{title} 欄位，可直接用多語言 SBERT 模型
重抽 384 維文字特徵，替換目前的 50 維舊版本：
\begin{lstlisting}[language=bash]
.venv/bin/python scripts/build_hetero_graph.py --use-sbert
.venv/bin/python scripts/train_hetero.py --model sage --epochs 300
\end{lstlisting}
這不需要下載任何新資料，首次執行會自動下載 SBERT 模型（約 400 MB）。
\end{tipbox}

% ════════════════════════════════════════════════════
\section{系統設計：三條平行管線}
\label{sec:system}

本報告同時維護三條研究管線，各自解決不同的問題。

\begin{figure}[H]
\centering
\begin{tikzpicture}[
  pipe/.style  ={draw=gray!60, fill=lightbg, rounded corners=8pt,
                 minimum width=4.0cm, minimum height=8.5cm,
                 text width=3.7cm, align=center},
  ptitle/.style={font=\bfseries\small, text=white, rounded corners=4pt,
                 minimum width=3.6cm, minimum height=0.8cm, align=center},
  pitem/.style ={font=\footnotesize, text=darktext, align=left, text width=3.4cm}]

\node[pipe] (pA) at (0,0) {};
\node[ptitle, fill=vidblue!80!black] at (0, 3.5) {Pipeline A};
\node[pitem] at (0, 2.6) {\textbf{問題}：哪支影片推給誰？};
\node[pitem] at (0, 1.7) {\textbf{方法}：MMGCN\\（多模態圖推薦）};
\node[pitem] at (0, 0.8) {\textbf{需要}：使用者歷史互動};
\node[pitem] at (0,-0.2) {\textbf{圖型}：使用者-影片二部圖};
\node[pitem] at (0,-1.1) {\textbf{輸出}：個人化影片排名};
\node[pitem] at (0,-2.2) {\textcolor{red!60!black}{冷啟動：無法處理}};

\node[pipe] (pB) at (5,0) {};
\node[ptitle, fill=authgreen!80!black] at (5, 3.5) {Pipeline B};
\node[pitem] at (5, 2.6) {\textbf{問題}：影片平均表現如何？};
\node[pitem] at (5, 1.7) {\textbf{方法}：LightGBM\\（梯度提升樹）};
\node[pitem] at (5, 0.8) {\textbf{需要}：影片特徵向量};
\node[pitem] at (5,-0.2) {\textbf{圖型}：無圖（純表格）};
\node[pitem] at (5,-1.1) {\textbf{輸出}：8 個互動指標};
\node[pitem] at (5,-2.2) {\textcolor{hlgold!80!black}{冷啟動：勉強可用}};

\node[pipe] (pC) at (10,0) {};
\node[ptitle, fill=catred!80!black] at (10, 3.5) {Pipeline C（新）};
\node[pitem] at (10, 2.6) {\textbf{問題}：新影片表現如何？};
\node[pitem] at (10, 1.7) {\textbf{方法}：HeteroGNN\\（異質圖神經網路）};
\node[pitem] at (10, 0.8) {\textbf{需要}：影片內容 + 作者 + 分類};
\node[pitem] at (10,-0.2) {\textbf{圖型}：影片中心異質圖};
\node[pitem] at (10,-1.1) {\textbf{輸出}：8 個互動指標};
\node[pitem] at (10,-2.2) {\textcolor{authgreen!80!black}{冷啟動：主要目標}};

\draw[-Stealth, thick, gray!50] (pA.east) -- (pB.west)
  node[midway, above, font=\scriptsize, text=gray] {個人→整體};
\draw[-Stealth, thick, gray!50] (pB.east) -- (pC.west)
  node[midway, above, font=\scriptsize, text=gray] {無圖→有圖};
\end{tikzpicture}
\caption{三條管線的設計定位。Pipeline A 解決個人化推薦，Pipeline B 提供無圖基線，
Pipeline C 是核心創新，用異質圖在冷啟動場景下預測影片表現。}
\label{fig:pipeline_overview}
\end{figure}

% ════════════════════════════════════════════════════
\section{Pipeline C：為什麼要建圖？}
\label{sec:why_graph}

\subsection{LightGBM（Pipeline B）的根本限制}

LightGBM 是一棵決策樹的集合，把每支影片的特徵向量（819 個數字）直接對應到
預測值。它無法知道「這支影片的作者上個月發的其他影片都爆紅了」這類\textbf{關係性信息}，
因為這個信息不在 819 個數字裡。

\subsection{圖的力量：借用鄰居的信息}

\begin{analogybox}[為什麼建圖有用？——用找房子類比]
你要估計一間剛上市、完全沒有租賃記錄的公寓的市場表現。
如果你只看這間公寓本身（幾坪、幾樓、幾個房間），你能得到的信息有限。

但如果你還能看到：
\begin{itemize}
\item 這個房東名下的其他公寓出租狀況都怎樣（作者節點的信號）
\item 同一個社區的其他公寓平均租出幾天（分類節點的信號）
\item 和這間公寓格局最像的 10 間公寓市場表現（相似影片節點的信號）
\end{itemize}
你的估計就會準確多了。GNN 的圖結構就是把這三種「鄰居信息」通道建立起來，
讓模型可以在訊息傳遞時自動學習如何加權使用這三種信號。
\end{analogybox}

% ════════════════════════════════════════════════════
\section{Pipeline C：圖的建構}
\label{sec:graph}

\subsection{節點設計：圖中有什麼？}

我們的圖有三種類型的節點，每種節點代表不同的實體：

\begin{table}[H]
\centering
\caption{三種節點類型、數量與特徵設計}
\begin{tabular}{llrl}
\toprule
\textbf{節點類型} & \textbf{代表} & \textbf{數量} & \textbf{特徵說明} \\
\midrule
影片（video）   & 每支影片 & 54,088 & 視覺 768 + 文字 50 + 時長 1 = \textbf{819 維} \\
作者（author）  & 每位作者 & 34,939 & 粉絲數 + 作品數 + 歷史均值 = \textbf{10 維} \\
分類（category）& 每個分類 & 51     & 可學習嵌入向量（\textbf{128 維}，隨訓練更新） \\
\bottomrule
\end{tabular}
\end{table}

\textbf{作者特徵的特殊設計（防資料洩漏）}：
作者的「8 指標歷史均值」只能用訓練集的影片統計，絕對不可以包含驗證集或測試集的數據。
若違反此規則，模型在評估時就等於「看到了答案」，指標會虛高。
程式碼中有明確的 assert 確保此規則：

\begin{lstlisting}[language=python]
# build_hetero_graph.py 中的防洩漏檢查
def build_author_features(meta_df, vid_attrs, train_mask):
    train_pids = set(meta_df.loc[train_mask, "pid"])
    assert used_pids.issubset(train_pids), "洩漏！"
\end{lstlisting}

\subsection{邊的設計：節點之間如何連接？}

\begin{figure}[H]
\centering
\begin{tikzpicture}[node distance=2.0cm and 2.5cm, scale=0.9]
\node[vidnode] (v1) at (0, 0)    {影片 $v_1$};
\node[vidnode] (v2) at (0, 2.2)  {影片 $v_2$};
\node[vidnode] (v3) at (0,-2.2)  {影片 $v_3$};
\node[vidnode, draw=vidblue!60, fill=vidblue!30, very thick]
               (vnew) at (-3.5, 0) {{\bfseries 新影片} $v_?$};
\node[authnode] (a1) at (5.5, 2.0)  {作者 $a_1$};
\node[authnode] (a2) at (5.5,-0.5)  {作者 $a_2$};
\node[catnode]  (c1) at (5.0,-2.8)  {分類 $c_1$};

\draw[-Stealth, draw=authgreen, thick]
  (v1) -- (a1) node[midway, above, sloped, font=\scriptsize, text=authgreen]
  {posted\_by（誰拍的）};
\draw[-Stealth, draw=authgreen, thick] (v2) -- (a1);
\draw[-Stealth, draw=authgreen, thick]
  (v3) -- (a2) node[midway, below, sloped, font=\scriptsize, text=authgreen]
  {posted\_by};
\draw[-Stealth, draw=catred, thick]
  (v1) -- (c1) node[midway, below, sloped, font=\scriptsize, text=catred]
  {belongs\_to（哪個分類）};
\draw[-Stealth, draw=catred, thick] (v3) -- (c1);
\draw[<->, dashed, draw=vidblue!70, thick]
  (v1) -- (v2) node[midway, right, font=\scriptsize, text=vidblue]
  {similar\_to（視覺相似）};
\draw[<->, dashed, draw=vidblue!70, thick] (v1) -- (v3);
\draw[-Stealth, draw=vidblue!60, thick, densely dashed]
  (vnew) -- (v1) node[midway, above, sloped, font=\scriptsize, text=vidblue]
  {cosine kNN};
\draw[-Stealth, draw=vidblue!60, thick, densely dashed]
  (vnew) to[bend left=15] (v2);

\begin{scope}[shift={(0,-4.8)}]
\draw[-, draw=authgreen, thick] (0,0)--(0.8,0) node[right, font=\scriptsize] {posted\_by：影片$\rightarrow$作者};
\draw[-, draw=catred, thick]    (4,0)--(4.8,0) node[right, font=\scriptsize] {belongs\_to：影片$\rightarrow$分類};
\draw[<->, dashed, draw=vidblue!70, thick] (8,0)--(8.8,0)
  node[right, font=\scriptsize] {similar\_to：視覺 cosine kNN};
\end{scope}
\end{tikzpicture}
\caption{異質圖結構。藍色節點是影片，綠色方框是作者，紅色菱形是分類。
加粗藍框是待預測的全新影片，它透過 cosine 最近鄰連入圖中——即使沒有互動記錄，
也能從鄰居那裡「借用」信號。}
\label{fig:hetero_graph}
\end{figure}

\textbf{cosine kNN 邊的建立方式}：
找每支影片視覺特徵最相似的 10 支影片，建立雙向邊。
關鍵限制：\textbf{只用視覺特徵計算相似度}，絕對不用任何行為標籤或互動資料，
確保這條邊在冷啟動時（新影片也有視覺特徵）同樣可用。

相似度計算公式（余弦相似度）：

$$\text{cosine}(v, v') = \frac{\mathbf{f}_v \cdot \mathbf{f}_{v'}}{\|\mathbf{f}_v\|\,\|\mathbf{f}_{v'}\|}$$

\textbf{白話翻譯}：把每支影片的 768 維特徵想成一個指向某方向的箭頭，
兩個箭頭越指向同一方向，cosine 值越接近 1，代表視覺風格越相似。
相似度取前 10 名的影片，連成 kNN 邊。

\subsection{圖的規模}

\begin{table}[H]
\centering
\caption{建立後的異質圖統計}
\begin{tabular}{lrr}
\toprule
\textbf{項目} & \textbf{數量} & \textbf{說明} \\
\midrule
影片節點 & 54,088 & 含訓練/驗證/測試集 \\
作者節點 & 34,939 & \\
分類節點 & 51     & 三級分類體系 \\
\midrule
posted\_by 邊（含反向） & 108,176 & 每部影片$\leftrightarrow$作者 \\
belongs\_to 邊（含反向）& 108,176 & 每部影片$\leftrightarrow$分類 \\
similar\_to 邊（雙向）  & 942,894 & $k=10$ cosine kNN \\
\textbf{邊的總數} & \textbf{1,159,246} & \\
\bottomrule
\end{tabular}
\end{table}

% ════════════════════════════════════════════════════
\section{Pipeline C：模型架構}
\label{sec:model}

\subsection{整體設計思路}

整個模型分成三個階段：
\begin{enumerate}
\item \textbf{輸入投影}：把三種不同維度的節點特徵，統一投影到同一個 128 維的空間
\item \textbf{訊息傳遞}：跑 2 層異質圖卷積，每層讓每個節點從鄰居那裡收集信息
\item \textbf{輸出頭}：只取影片節點的 128 維表示，接一個多層感知機，輸出 8 個指標
\end{enumerate}

\begin{figure}[H]
\centering
\begin{tikzpicture}[node distance=0.55cm]
\node[layerbox] (input) {\textbf{輸入層}\\
  \footnotesize video.x $[54088 \times 819]$，author.x $[34939 \times 10]$，category embedding};
\node[layerbox, below=of input] (proj)
  {\textbf{輸入投影}（統一維度）\\
  \footnotesize Linear($-1$, 128) 對影片/作者；Embedding(51, 128) 對分類\\
  $\Rightarrow$ 所有節點都變成 128 維};
\node[layerbox, below=of proj, draw=vidblue!60, fill=vidblue!8] (conv1)
  {\textbf{HeteroConv 第 1 層}（1 跳訊息傳遞）\\
  \footnotesize 每個節點收集所有直接鄰居的 128 維向量，加總後更新自己\\
  $\Rightarrow$ 影片「知道」作者、分類、相似影片的原始特徵；ReLU + Dropout};
\node[layerbox, below=of conv1, draw=vidblue!60, fill=vidblue!8] (conv2)
  {\textbf{HeteroConv 第 2 層}（2 跳訊息傳遞）\\
  \footnotesize 再收集一次——現在鄰居的向量已包含了它的鄰居的信息\\
  $\Rightarrow$ 影片的感受野擴展到 2 跳以內的所有節點；ReLU + Dropout};
\node[layerbox, below=of conv2, draw=catred!60, fill=catred!8] (head)
  {\textbf{輸出頭（只取影片節點）}\\
  \footnotesize $\mathbf{h}_v^{(2)}$ [128 維] → Linear(128,64) → ReLU → Linear(64,8)};

\draw[arrowstyle] (input) -- (proj);
\draw[arrowstyle] (proj)  -- (conv1);
\draw[arrowstyle] (conv1) -- (conv2);
\draw[arrowstyle] (conv2) -- (head);

\node[draw=authgreen, fill=authgreen!10, rounded corners=3pt,
      font=\scriptsize, right=0.3cm of proj, text width=3cm, align=left]
  {參數總量：\\450,632 個\\（在 RTX 4090 上\\約 2 MB 顯存）};
\end{tikzpicture}
\caption{HeteroGNN 完整模型架構。輸入投影將各種節點統一到 128 維，
兩層 HeteroConv 逐步擴大感受野，最後只在影片節點上接多任務輸出頭。}
\label{fig:arch}
\end{figure}

\subsection{訊息傳遞的數學表達}

一層異質圖卷積的更新公式：

\begin{equation}
\mathbf{h}_v^{(\ell+1)} = \sigma\!\Biggl(\;
\sum_{\text{邊類型}\; r} \;\mathbf{W}_r^{(\ell)} \cdot
\underbrace{\frac{1}{|\mathcal{N}_r(v)|}
  \sum_{u \in \mathcal{N}_r(v)} \mathbf{h}_u^{(\ell)}}_{\text{關係 $r$ 下的鄰域平均}}
\;\Biggr)
\label{eq:hetero_msg}
\end{equation}

\textbf{每個符號的白話解釋}：
\begin{itemize}
\item $\mathbf{h}_v^{(\ell)}$：第 $\ell$ 層後，節點 $v$（某支影片）的 128 維特徵向量
\item $r$：邊的類型（posted\_by / belongs\_to / similar\_to 及其反向邊）
\item $\mathcal{N}_r(v)$：在關係 $r$ 下，$v$ 的鄰居集合（例如 $r = \text{posted\_by}$ 時，$\mathcal{N}_r(v)$ 就是這支影片的作者）
\item $\mathbf{W}_r^{(\ell)}$：關係 $r$ 的可學習投影矩陣——模型訓練時自動學習「作者信號」比「分類信號」在哪些情況下更重要
\item $\sigma$：ReLU 激活函數（讓模型學到非線性關係）
\end{itemize}

\textbf{整句白話翻譯}：「對影片 $v$，對每種鄰居關係，把那種關係下的所有鄰居特徵平均，然後用該關係專屬的投影矩陣轉換，最後把所有關係的結果加起來，通過 ReLU，得到下一層的特徵。」

\begin{figure}[H]
\centering
\begin{tikzpicture}[scale=0.90]
\node[vidnode, draw=vidblue!30, fill=gray!10] (s1) at (0, 3.8) {相似影片 $v_1$};
\node[vidnode, draw=vidblue!30, fill=gray!10] (s2) at (0, 2.2) {相似影片 $v_2$};
\node[authnode, draw=authgreen!40, fill=gray!10] (au) at (0, 0.5) {作者 $a$};
\node[catnode,  draw=catred!40,   fill=gray!10] (ca) at (0,-1.2) {分類 $c$};

\node[font=\scriptsize, right=0.05cm of s1, text=gray] {h$^{(0)}$：視覺特徵};
\node[font=\scriptsize, right=0.05cm of au, text=gray] {h$^{(0)}$：作者歷史均值};
\node[font=\scriptsize, right=0.05cm of ca, text=gray] {h$^{(0)}$：分類嵌入};

\node[vidnode, very thick, fill=vidblue!25] (v1) at (5.5, 1.5)
  {新影片 $v_?$\\\small $\mathbf{h}^{(1)}$};

\draw[-Stealth, thick, draw=vidblue!70, dashed] (s1.east) -- (v1.west)
  node[midway, above, sloped, font=\scriptsize, text=vidblue] {similar\_to};
\draw[-Stealth, thick, draw=vidblue!70, dashed] (s2.east) -- (v1.west);
\draw[-Stealth, thick, draw=authgreen] (au.east) -- (v1.west)
  node[midway, below, sloped, font=\scriptsize, text=authgreen] {posted\_by};
\draw[-Stealth, thick, draw=catred] (ca.east) -- (v1.west)
  node[midway, below, sloped, font=\scriptsize, text=catred] {belongs\_to};

\node[vidnode, very thick, fill=vidblue!45] (v2) at (5.5,-0.5)
  {新影片 $v_?$\\\small $\mathbf{h}^{(2)}$（2 跳）};
\draw[arrowstyle, very thick] (v1) -- (v2)
  node[midway, right, font=\scriptsize, text=gray] {Layer 2\\再聚合一次};

\node[draw=hlgold!80, fill=hlgold!15, rounded corners=3pt,
      font=\small, align=center, minimum width=2.5cm] (out) at (5.5,-2.5)
  {MLP → 8 個指標};
\draw[arrowstyle] (v2) -- (out);

\node[draw=gray!30, fill=gray!5, rounded corners=3pt,
      text width=4.5cm, font=\scriptsize, align=left] at (10, 1.2)
  {\textbf{冷啟動場景：}\\[3pt]
   新影片沒有任何互動記錄，\\但它依然能收集信號：\\[3pt]
   \textcolor{vidblue}{$\bullet$ 相似影片}的歷史表現\\
   \textcolor{authgreen}{$\bullet$ 作者}的過去作品均值\\
   \textcolor{catred}{$\bullet$ 分類}的整體基準\\[3pt]
   不需要任何使用者互動。};
\end{tikzpicture}
\caption{冷啟動場景的訊息傳遞示意。全新影片 $v_?$ 沒有互動記錄，但透過三種邊
從圖的鄰居收集信號，兩層後已融合了 2 跳鄰域的資訊，最後通過 MLP 輸出預測。}
\label{fig:msgpass}
\end{figure}

\subsection{兩種訊息傳遞變體}

我們提供兩種可切換的訊息傳遞層：

\begin{table}[H]
\centering
\caption{GraphSAGE（預設）vs HGT（消融用）}
\begin{tabular}{lll}
\toprule
\textbf{面向} & \textbf{GraphSAGE（\texttt{-{}-model sage}）} & \textbf{HGT（\texttt{-{}-model hgt}）} \\
\midrule
聚合方式 & 鄰域平均（簡單直接） & 注意力加權（讓模型決定聽誰的）\\
直覺解釋 & 等量採納所有鄰居意見 & 根據內容決定聽哪個鄰居的話\\
泛化能力 & 強（推薦用於冷啟動） & 中等 \\
計算成本 & 低 & 高（4 個注意力頭）\\
\bottomrule
\end{tabular}
\end{table}

\textbf{注意力機制（HGT 的核心）白話解釋}：
普通 GNN 對所有鄰居一視同仁，把它們的特徵平均起來。
加了注意力機制後，模型會學習「對這支影片來說，作者的信號比分類的信號更重要」，
對每個鄰居給予不同的權重，讓不同情況下的信息融合更有彈性。

% ════════════════════════════════════════════════════
\section{訓練設計}
\label{sec:training}

\subsection{為什麼要「加權」訓練？}

同樣是「按讚率 5\%」，一支被 1000 人看過的影片的統計比一支只被 10 人看過的影片要可靠多了。
後者可能純屬巧合。因此訓練時，每支影片的貢獻正比於它的曝光次數：

$$w_v = \frac{n_v}{\sum_{u \in \text{訓練集}} n_u} \qquad (\text{n\_impressions 歸一化})$$

\subsection{多任務損失函數}

8 個目標同時訓練，損失函數分兩部分：

\begin{equation}
\mathcal{L} = \sum_{v \in \text{訓練集}} w_v \Biggl[
  \underbrace{\sum_{k=0}^{6}\bigl(\hat{y}_{v,k} - y_{v,k}\bigr)^2}_{\text{7 個 rate 指標用 MSE}}
  \;+\;
  \underbrace{\text{Huber}\!\bigl(\hat{y}_{v,7},\,\log(1+y_{v,7});\,\delta{=}1.0\bigr)}_{\text{觀看時長用 Huber（更穩健）}}
\Biggr]
\end{equation}

\textbf{為什麼觀看時長用 Huber 損失而不是 MSE？}

觀看時長的分佈非常偏斜：大部分人看不到 30 秒，但有少數影片被看了好幾分鐘。
MSE 對「極端值」的懲罰非常大（誤差平方），一兩個極端樣本就會主導整個訓練。
Huber 損失在誤差小的時候用平方（精確），在誤差大的時候改用線性（不被極端值帶跑），
兼顧了穩定性和準確性。

\textbf{為什麼觀看時長要先取對數？}
同樣的道理：直接對秒數訓練，500 秒和 5 秒的差距（495 秒）會把差距很小的
5 秒和 10 秒（5 秒）完全淹沒。取 $\log(1 + t)$ 後，這個比例差異被壓縮成相近的數值，
讓模型對各種時長範圍都能均等地學習。

\subsection{資料切分與防洩漏設計}

\begin{figure}[H]
\centering
\begin{tikzpicture}[scale=0.9]
\fill[vidblue!20, rounded corners=4pt] (0,0) rectangle (10.5, 1.2);
\node[font=\small] at (5.25, 0.6) {全部 54,088 部影片（工作集）};
\fill[vidblue!50, rounded corners=3pt] (0.1, 0.1) rectangle (8.4, 1.1);
\node[font=\footnotesize, text=white] at (4.25, 0.6) {訓練集 43,270（80\%）};
\fill[authgreen!50, rounded corners=3pt] (8.5, 0.1) rectangle (9.45, 1.1);
\node[font=\scriptsize, text=white] at (8.97, 0.6) {驗\\5,408};
\fill[catred!50, rounded corners=3pt] (9.55, 0.1) rectangle (10.4, 1.1);
\node[font=\scriptsize, text=white] at (9.97, 0.6) {測\\5,410};
\draw[red!70!black, very thick, dashed] (8.45,-0.4) -- (8.45, 3.0)
  node[above, font=\small, text=red!70!black] {標籤洩漏屏障};
\node[draw=authgreen, fill=authgreen!10, rounded corners=4pt,
      text width=5cm, align=center, font=\small] (astat) at (4.25, 2.2)
  {作者歷史統計\\僅用訓練集影片計算};
\draw[-Stealth, authgreen, thick] (4.25, 1.2) -- (astat.south);
\node[font=\Large, text=red!70!black] at (9.2, 2.2) {$\times$};
\node[font=\small, text=red!70!black, align=center] at (9.2, 1.8) {禁止};
\node[draw=gray!40, fill=gray!5, rounded corners=3pt,
      text width=4cm, font=\scriptsize, align=center] at (4.25,-1.2)
  {程式碼 assert 確保覆蓋率 = 100\%\\建圖時印出確認訊息};
\end{tikzpicture}
\caption{80/10/10 資料切分與防洩漏屏障。作者的歷史統計特徵嚴格限定在訓練集範圍內，
程式碼中以 assert 明確驗證，確保評估時不存在資訊洩漏。}
\label{fig:leakage}
\end{figure}

% ════════════════════════════════════════════════════
\section{實驗結果}
\label{sec:exp}

\subsection{實驗設定}

\begin{itemize}
\item 資料切分：80/10/10，seed=42（與 Pipeline B LightGBM 完全一致，保證公平比較）
\item 訓練環境：Ubuntu 24.04，NVIDIA RTX 4090（24 GB VRAM），PyTorch 2.6.0+cu124
\item 早停：監控驗證集損失，patience=50
\item 最終評估：使用最佳驗證 checkpoint 在測試集評估
\end{itemize}

\subsection{GPU 訓練加速效果}

\begin{table}[H]
\centering
\caption{CPU 與 GPU 訓練結果對比（seed=42，SAGE 模型）}
\begin{tabular}{lrrrr}
\toprule
\textbf{環境} & \textbf{早停 epoch} & \textbf{訓練時間} & \multicolumn{2}{c}{\textbf{mean\_watch\_time}} \\
\cmidrule(lr){4-5}
& & & Spearman & MAE（秒）\\
\midrule
macOS CPU（Apple M 系列） & 66 & 47 秒 & 0.251 & 505 \\
Ubuntu RTX 4090 & 68 & \textbf{1 秒} & \textbf{0.256} & \textbf{386} \\
\midrule
加速比 & — & \textbf{47×} & $\uparrow$ 2\% & $\uparrow$ 24\% \\
\bottomrule
\end{tabular}
\end{table}

\subsection{主要指標對比（GPU 訓練結果）}

\begin{table}[H]
\centering
\caption{Pipeline C HeteroGNN vs Pipeline B LightGBM（測試集，GPU 結果）}
\label{tab:results}
\begin{tabular}{lrrrr}
\toprule
& \multicolumn{2}{c}{\textbf{Spearman（越高越好）}}
& \multicolumn{2}{c}{\textbf{MAE（越低越好）}} \\
\cmidrule(lr){2-3}\cmidrule(lr){4-5}
\textbf{指標} & \textbf{HeteroGNN} & \textbf{LightGBM}
              & \textbf{HeteroGNN} & \textbf{LightGBM} \\
\midrule
\texttt{like\_rate}            & $-0.032$ & $-0.001$ & $0.058$ & $0.044$ \\
\texttt{comment\_rate}         & $+0.012$ & $-0.003$ & $0.039$ & $0.009$ \\
\texttt{follow\_rate}          & $-0.031$ & $-0.009$ & $0.024$ & $0.008$ \\
\texttt{collect\_rate}         & $+0.011$ & $-0.035$ & $0.074$ & $0.013$ \\
\texttt{forward\_rate}         & $-0.005$ & $-0.012$ & $0.066$ & $0.005$ \\
\texttt{hate\_rate}            & $+0.019$ & $-0.014$ & $0.029$ & $0.002$ \\
\texttt{effective\_view\_rate} & $-0.012$ & $+0.016$ & $0.169$ & $0.143$ \\
\texttt{mean\_watch\_time}     & $\mathbf{+0.256}$ & $+0.047$ & $386$ 秒 & $24$ 秒 \\
\bottomrule
\end{tabular}
\end{table}

\subsection{結果解讀}

\begin{tipbox}[最重要的發現]
\textbf{mean\_watch\_time 的排名能力提升 4 倍以上}：Spearman $0.047 \to 0.256$。
這驗證了本報告的核心假設：圖結構確實能有效傳遞「相似影片的觀看時長信號」。
\end{tipbox}

\paragraph{觀看時長：圖結構有效}
觀看時長是最受影片「內容本身」影響的指標——有趣的影片讓你多看幾秒，無聊的讓你提前划走。
GNN 透過 \texttt{similar\_to} 邊，把「視覺風格相似的老影片的觀看時長」傳給新影片，
這個信號是 LightGBM 看不到的。

\paragraph{Rate 指標：兩個模型都接近隨機}
按讚、留言、轉發等 rate 的中位數是 0，意味著絕大多數影片根本沒人互動。
「哪支影片會爆紅」有很強的隨機性和社群動態（例如一個大號轉發就能帶動流量），
純從影片內容預測這種「社群爆紅效應」是學術界公認的困難問題。
兩個模型的 Spearman 接近 0 是正常且預期的結果。

\paragraph{MAE 的矛盾：排名好但絕對值差}
HeteroGNN 的觀看時長 MAE（386 秒）遠高於 LightGBM（24 秒），
但 Spearman（0.256）遠優於 LightGBM（0.047）。
這意味著：GNN 能夠正確排序「哪支影片會被看更久」，但絕對的預測秒數還需要更多訓練和校準。
對推薦系統應用來說，\textbf{排名能力（Spearman）比絕對精度（MAE）更重要}。

% ════════════════════════════════════════════════════
\section{未來工作}
\label{sec:future}

按優先順序排列：

\subsection*{高優先：補齊功能缺口}

\paragraph{1. 冷啟動推論的精確節點對應}
當前 \texttt{predict\_hetero.py from-raw} 模式若提供 \texttt{--author-id}，
因建圖時的對應字典沒有持久化，只能退回 kNN fallback。
需額外儲存 \texttt{hetero\_graph\_meta.pkl} 包含作者和分類的本地索引字典。

\paragraph{2. Category 三層階層嵌入}
圖中已存 \texttt{category.hier}（含 leaf/parent/root 三層索引），
但模型只用了單一嵌入。應改為三層各自一個 Embedding 後拼接，
充分利用「這個分類$\rightarrow$它的父分類$\rightarrow$根分類」的語義層次。

\subsection*{中優先：提升模型能力}

\paragraph{3. SBERT 文字特徵}（已支援 \texttt{--use-sbert}）
把文字特徵從 50 維提升至 384 維，特別有利於標題語義豐富的影片。

\paragraph{4. 多任務不確定性加權}（Kendall 2018）
目前 8 個任務等權重，改用可學習的 $\sigma_k^2$ 自動平衡：
$\mathcal{L} = \sum_k \frac{1}{2\sigma_k^2} \mathcal{L}_k + \log \sigma_k$

\paragraph{5. Tag 節點}
加入 $(\text{video}, \texttt{has\_tag}, \text{tag})$ 邊，
利用 \texttt{interaction.csv} 中的 \texttt{tag\_name} 語義信號。

\subsection*{低優先：評估完整性}

\begin{itemize}
\item 多 seed 評估（seed=0/1/2/42），獲得統計顯著的均值±標準差
\item HGT vs SAGE 消融實驗
\item 與 MMVED/HMMVED 公開 baseline 在同一測試集比較
\item 真正的歸納冷啟動測試（held-out 的完全未見影片集）
\end{itemize}

% ════════════════════════════════════════════════════
\section{相關工作}
\label{sec:related}

\paragraph{圖神經網路基礎}
GCN \citep{kipf2017semi}奠定了訊息傳遞框架；
GraphSAGE \citep{hamilton2017inductive}解決了歸納推論問題（能處理訓練時未見過的新節點，
正是冷啟動的必要條件）；
GAT \citep{velivckovic2018graph}引入注意力機制讓鄰居貢獻可學習加權；
HGT \citep{hu2020heterogeneous}將注意力推廣至異質圖，是本報告 HGT 變體的直接來源。

\paragraph{影片人氣預測}
MMVED \citep{xie2020mmved}及其後續 HMMVED \citep{xie2021hmmved}是最直接的對標 baseline，
以多模態融合從影片內容預測人氣，有公開程式碼。
GraphInf \citep{graphinf2022}驗證了圖卷積在短影片人氣預測上的有效性。

\paragraph{多任務不確定性加權}
\citet{kendall2018multi}提出以可學習的任務不確定性動態平衡多任務損失，
是本報告後續改進的主要候補方案。

% ════════════════════════════════════════════════════
\section{結論}

本報告完成了三項核心貢獻：
\begin{enumerate}
\item \textbf{框架設計}：建立了三管線並行研究框架，清晰區分推薦（Pipeline A）、
非圖基線（Pipeline B）、圖基線（Pipeline C）三個不同問題設定。
\item \textbf{技術驗證}：Pipeline C 在 RTX 4090 GPU 上的實驗表明，
異質圖結構在觀看時長排名預測上有效（Spearman $0.047 \to 0.256$，提升 4 倍以上）。
\item \textbf{工程實作}：完整實作了從建圖、訓練、到冷啟動推論的端對端 pipeline，
並加入嚴格的防資料洩漏機制。
\end{enumerate}

下一步將重點在 GPU 環境上展開大模型實驗（$d=256$，3 層），
並引入 SBERT 文字特徵與多任務不確定性加權，
目標是讓 \texttt{mean\_watch\_time} Spearman 突破 0.35。

% ════════════════════════════════════════════════════
\bibliographystyle{abbrvnat}
\begin{thebibliography}{99}

\bibitem[Hamilton et al.(2017)]{hamilton2017inductive}
Hamilton, W., Ying, Z., and Leskovec, J.
\newblock \emph{Inductive Representation Learning on Large Graphs}.
\newblock \textit{NeurIPS}, 2017.

\bibitem[Hu et al.(2020)]{hu2020heterogeneous}
Hu, Z., Dong, Y., Wang, K., and Sun, Y.
\newblock \emph{Heterogeneous Graph Transformer}.
\newblock \textit{WWW}, 2020.

\bibitem[Kendall et al.(2018)]{kendall2018multi}
Kendall, A., Gal, Y., and Cipolla, R.
\newblock \emph{Multi-Task Learning Using Uncertainty to Weigh Losses}.
\newblock \textit{CVPR}, 2018.

\bibitem[Kipf and Welling(2017)]{kipf2017semi}
Kipf, T.~N. and Welling, M.
\newblock \emph{Semi-Supervised Classification with Graph Convolutional Networks}.
\newblock \textit{ICLR}, 2017.

\bibitem[GraphInf(2022)]{graphinf2022}
GraphInf Authors.
\newblock \emph{Graph-Based Content Popularity Prediction for Short Videos}.
\newblock Technical Report, 2022.

\bibitem[Shang et al.(2025)]{shang2025shortvideo}
Shang, J. et al.
\newblock \emph{A Large-scale Dataset with Behavior, Attributes, and Content
  of Mobile Short-video Platform}.
\newblock \textit{WWW}, 2025.

\bibitem[Velickovic et al.(2018)]{velivckovic2018graph}
Veli\v{c}kovi\'{c}, P. et al.
\newblock \emph{Graph Attention Networks}.
\newblock \textit{ICLR}, 2018.

\bibitem[Wei et al.(2019)]{wei2019mmgcn}
Wei, Y. et al.
\newblock \emph{MMGCN: Multi-modal Graph Convolution Network for Personalized
  Recommendation of Micro-video}.
\newblock \textit{ACM MM}, 2019.

\bibitem[Xie et al.(2020)]{xie2020mmved}
Xie, H. et al.
\newblock \emph{MMVED: Multi-modal Variational Encoder-Decoder for Micro-video
  Popularity Prediction}.
\newblock \textit{WWW}, 2020.

\bibitem[Xie et al.(2021)]{xie2021hmmved}
Xie, H. et al.
\newblock \emph{A Hierarchical Multi-task Learning Framework for Micro-video
  Popularity Prediction}.
\newblock \textit{IEEE TMM}, 2021.

\end{thebibliography}

\end{document}
